%! Matrices Modeling Contexts — Unit 4, Lesson 4.14 — Warm-Up Key
\documentclass[10pt]{article}
\usepackage{precalculus-article}
\usepackage{precalculus-key}
\newcommand{\TallMath}[1]{$\displaystyle #1\rule[-1.4em]{0pt}{3.2em}$}

\begin{document}

\pageheader{Unit 4, Lesson 4.14}{Warm-Up Key}
\namedateperiod

\vspace{0.15in}

\begin{notesbox}{Warm-Up Key}

\textbf{1.}\quad Compute the matrix--vector product. Show your work in the space provided.

\[
\begin{bmatrix}0.7 & 0.4 \\ 0.3 & 0.6\end{bmatrix}
\begin{bmatrix}100 \\ 80\end{bmatrix}
= \begin{bmatrix}\ans{102} \\ \ans{78}\end{bmatrix}
\]

\vspace{10pt}
\textit{Work:}\quad $0.7(100)+0.4(80) = \ans{70+32=102}$

\vspace{6pt}
\phantom{\textit{Work:}\quad }$0.3(100)+0.6(80) = \ans{30+48=78}$

\vspace{18pt}
\textbf{2.}\quad Given $A = \begin{bmatrix}2&1\\0&3\end{bmatrix}$ and
$A^{-1} = \begin{bmatrix}1/2 & -1/6 \\ 0 & 1/3\end{bmatrix}$,
verify that $A \cdot A^{-1} = I$ by computing the product.

\vspace{8pt}
$A \cdot A^{-1} = \begin{bmatrix}2&1\\0&3\end{bmatrix}\begin{bmatrix}1/2&-1/6\\0&1/3\end{bmatrix}
= \ans{$\begin{bmatrix}1&0\\0&1\end{bmatrix}$}$

\vspace{8pt}
\textit{Work (top row):}\quad $2(1/2)+1(0) = \ans{1}$ \qquad $2(-1/6)+1(1/3) = \ans{-1/3+1/3=0}$

\vspace{6pt}
\textit{Work (bottom row):}\quad $0(1/2)+3(0) = \ans{0}$ \qquad $0(-1/6)+3(1/3) = \ans{1}$

\vspace{18pt}
\textbf{3.}\quad A column of a transition matrix records the proportions leaving each state.
The two entries in any column must sum to \ans{1}. Why?

\vspace{4pt}
\ansline{Each column represents 100\% of the members of one state --- every person must go somewhere,}
\ansline{so the proportions must account for the entire group and sum to 1.}

\end{notesbox}

\end{document}
